% Figure: Load-Bearing versus Secondary Locks
% Structural analogy for disease-sustaining mechanisms in ME/CFS

\begin{figure}[htbp]
\centering
\begin{tikzpicture}[
    scale=0.80, every node/.style={scale=0.80, font=\sffamily},
    % Wall styles
    loadbearing/.style={
        draw=red!50!black, fill=red!25, line width=2.5pt,
        minimum height=#1, minimum width=0.55cm,
        pattern=north east lines, pattern color=red!40!black,
        inner sep=0pt
    },
    loadbearing/.default=3.2cm,
    secondary/.style={
        draw=gray!60, fill=gray!15, line width=1pt,
        minimum height=#1, minimum width=0.3cm,
        inner sep=0pt
    },
    secondary/.default=3.2cm,
    % Label styles
    walllabel/.style={font=\sffamily\scriptsize, align=center, text width=2.2cm},
    paneltitle/.style={font=\sffamily\bfseries\normalsize, anchor=south},
]

% ============================================================
% LEFT PANEL: Disease Architecture
% ============================================================
\begin{scope}[shift={(0,0)}]

    % Panel title
    \node[paneltitle] at (3.5, 5.8) {Disease Architecture};

    % Roof / "Disease State" slab
    \fill[red!8, draw=red!40!black, thick, rounded corners=2pt]
        (-0.3, 4.6) rectangle (7.3, 5.4);
    \node[font=\sffamily\bfseries\small, text=red!50!black] at (3.5, 5.0)
        {ME/CFS Disease State};

    % Floor
    \fill[gray!30, draw=gray!60, thick]
        (-0.3, 0.0) rectangle (7.3, 0.3);
    \node[font=\sffamily\tiny, text=gray!60] at (3.5, 0.15) {baseline physiology};

    % --- Load-bearing walls (dark red, hatched, thick) ---

    % Wall 1: Epigenetic Consolidation (left)
    \node[loadbearing=4.0cm] (lb1) at (1.0, 2.5) {};
    \node[walllabel, below=0.15cm of lb1, text=red!50!black]
        {Epigenetic\\Consolidation};

    % Wall 2: Autoimmune Persistence (right)
    \node[loadbearing=4.0cm] (lb2) at (6.0, 2.5) {};
    \node[walllabel, below=0.15cm of lb2, text=red!50!black]
        {Autoimmune\\Persistence\\{\tiny(plasma cells)}};

    % --- Secondary walls (thin, light gray) ---

    % Wall 3: Oxidative Stress Cycle
    \node[secondary=3.6cm] (sw1) at (2.6, 2.3) {};
    \node[walllabel, below=0.15cm of sw1, text=gray!70]
        {Oxidative\\Stress Cycle};

    % Wall 4: Sleep Fragmentation
    \node[secondary=3.6cm] (sw2) at (3.8, 2.3) {};
    \node[walllabel, below=0.15cm of sw2, text=gray!70]
        {Sleep\\Fragmentation};

    % Wall 5: Viral Reactivation
    \node[secondary=3.6cm] (sw3) at (5.0, 2.3) {};
    \node[walllabel, below=0.15cm of sw3, text=gray!70]
        {Viral\\Reactivation};

    % Legend
    \node[loadbearing=0.4cm, anchor=west] (legLB) at (-0.1, -1.6) {};
    \node[font=\sffamily\scriptsize, anchor=west] at (0.5, -1.6)
        {Load-bearing lock};
    \node[secondary=0.4cm, anchor=west] (legSW) at (3.2, -1.6) {};
    \node[font=\sffamily\scriptsize, anchor=west] at (3.6, -1.6)
        {Secondary lock};

\end{scope}

% ============================================================
% RIGHT PANEL: Treatment Effect Prediction
% ============================================================
\begin{scope}[shift={(10, 0)}]

    % Panel title
    \node[paneltitle] at (2.8, 5.8) {Treatment Effect Prediction};

    % --- Top graph: Load-bearing lock removed (slope change) ---
    \begin{scope}[shift={(0, 3.0)}]
        % Axes
        \draw[-latex, thick, gray!80] (0, 0) -- (5.8, 0)
            node[font=\sffamily\scriptsize, below, pos=1] {Time};
        \draw[-latex, thick, gray!80] (0, 0) -- (0, 2.5)
            node[font=\sffamily\scriptsize, above, rotate=90, anchor=south, pos=0.5]
            {Severity};

        % Title
        \node[font=\sffamily\small\bfseries, text=red!50!black, anchor=west]
            at (0.2, 2.3) {Load-bearing lock removed};

        % Disease trajectory: high, then progressive decline after intervention
        \draw[very thick, red!40!black]
            (0.3, 1.8) -- (2.0, 1.8);
        \draw[very thick, red!50!black, -{Latex[length=2mm]}]
            (2.0, 1.8) -- (5.2, 0.4);

        % Intervention marker
        \draw[dashed, gray!60] (2.0, 0) -- (2.0, 2.0);
        \node[font=\sffamily\tiny, gray!60, anchor=south] at (2.0, 2.0) {Tx};

        % Annotation: slope change
        \draw[<->, thin, blue!60!black] (3.0, 1.8) -- (3.0, 1.27);
        \node[font=\sffamily\tiny, blue!60!black, anchor=west] at (3.1, 1.5)
            {slope $\Delta$};

        % Arrow indicating progressive improvement
        \node[font=\sffamily\tiny, text=red!50!black, anchor=west, text width=1.8cm]
            at (3.8, 0.6) {progressive\\recovery};
    \end{scope}

    % --- Bottom graph: Secondary lock removed (level change) ---
    \begin{scope}[shift={(0, -0.2)}]
        % Axes
        \draw[-latex, thick, gray!80] (0, 0) -- (5.8, 0)
            node[font=\sffamily\scriptsize, below, pos=1] {Time};
        \draw[-latex, thick, gray!80] (0, 0) -- (0, 2.5)
            node[font=\sffamily\scriptsize, above, rotate=90, anchor=south, pos=0.5]
            {Severity};

        % Title
        \node[font=\sffamily\small\bfseries, text=gray!60, anchor=west]
            at (0.2, 2.3) {Secondary lock removed};

        % Disease trajectory: high, then step down, then plateau
        \draw[very thick, gray!50]
            (0.3, 1.8) -- (2.0, 1.8);
        \draw[very thick, gray!50]
            (2.0, 1.8) -- (2.0, 1.1);
        \draw[very thick, gray!50, -{Latex[length=2mm]}]
            (2.0, 1.1) -- (5.2, 1.1);

        % Intervention marker
        \draw[dashed, gray!60] (2.0, 0) -- (2.0, 2.0);
        \node[font=\sffamily\tiny, gray!60, anchor=south] at (2.0, 2.0) {Tx};

        % Annotation: level change
        \draw[<->, thin, blue!60!black] (2.4, 1.8) -- (2.4, 1.1);
        \node[font=\sffamily\tiny, blue!60!black, anchor=west] at (2.5, 1.45)
            {level $\Delta$};

        % Annotation: same slope (flat)
        \node[font=\sffamily\tiny, text=gray!50, anchor=west, text width=1.8cm]
            at (3.5, 0.7) {symptomatic\\relief only};
    \end{scope}

\end{scope}

\end{tikzpicture}
\caption{Load-bearing versus secondary locks in ME/CFS. Load-bearing locks
    (epigenetic consolidation, autoimmune persistence) sustain the disease
    architecture; their removal enables progressive recovery (slope change).
    Secondary locks (oxidative stress, sleep fragmentation, viral reactivation)
    modulate severity; their removal provides symptomatic relief (level change)
    without altering trajectory (Section~\ref{sec:load-bearing-locks}).}
\label{fig:load-bearing-locks}
\end{figure}
